\newpage
\section{The Proposed System}

\subsection{Business Context}

\subsubsection{Application domain}

The application domain is career guidance and exploration, situated at the intersection of educational technology and occupational information provision. Participants in this domain are individuals facing an occupational decision, institutions supporting them, and the bodies that publish occupational information. The domain's characteristic activity is the transfer of occupational information to a person who must act on it, and Chapter 1 established that this transfer is currently incomplete: guidance provision is dominated by instruments that elicit self-description rather than convey what work involves.

The proposed system, JobQuest, addresses this by allowing a user to perform representative tasks of an occupation, receive an assessment of the competencies those tasks elicited, and use that assessment to determine which further occupations to explore. It is a platform for informed exploration rather than a placement, recruitment, or credentialing service, and this boundary determines several of the business rules below.

\subsubsection{Business processes}

Three processes constitute the operation of the platform. They are modelled formally in Section 5.1.1.

The exploration process is the primary user-facing process. A registered user optionally completes an orientation conversation, which is provided for users who do not know which occupation to begin with and is not a prerequisite. The user selects an occupation and a level within it, is presented with the role, team, and working context before any task begins, and then works through the tasks comprising that level. Each task is a scenario drawn from the occupation and is answered in one of nine formats --- multiple choice, multiple selection, ordering, matching, true/false, timed choice, free text, simplified email, or simplified code --- with hints available throughout. Tasks ask what the user does in the situation, a behavioural-tendency instruction\cite{r28}, because how a user tends to act is itself information about fit; scoring uses expert effectiveness keys in either case. Each response option is of one of three types: correct, the most effective response; trade-off, effective for one competency at a cost to another; or wrong, ineffective or harmful. Options and rubric criteria are scored on a single ordinal scale per competency, from $-$1 (ineffective or harmful) through 0 and +1 to +2 (most effective), so that evidence from choices and from written responses accumulates in the same competence record and agreement with experts can be computed on an ordinal scale. Using a hint is recorded and lowers the weight of that response as evidence of independent competence without blocking progress; a timed task that expires is recorded as unanswered, with its response time, and carries no separate penalty. The system evaluates each response, records the competencies it evidences, and returns a result, feedback specific to the response, and guidance for improvement. Completing a level awards competence credit; accumulated credit unlocks higher levels of the same occupation or, where competencies overlap sufficiently, entry to an adjacent occupation above its entry level. On completing an occupation, the user may rate it on five attributes, which contribute to that occupation's aggregate profile.

The content lifecycle process governs how occupations enter the catalogue, and it involves two roles with distinct responsibilities. Domain experts own the occupational content and author it in the system: for each scenario they enter the required facts, the competencies it must elicit, the critical incidents from which response options are drawn, the effectiveness of each option, and the rubric against which free responses are judged, which together form the scenario's seed. A seed is a template rather than a script: it fixes the canon events of a scenario --- the situations every run must pass through, at set points --- together with the keyed options and rubric attached to them, and leaves the narrative between those events to be developed by the model in response to the user's answers. Validation accordingly has two layers. In the first, the system generates a sample of scenarios from the seed, and experts other than the seed's authors review the seed and the sample in the system, judging their accuracy and returning approvals, rejections, or corrections, each recorded with the identity of the reviewing expert. In the second, every scenario generated at run time is checked automatically against its seed --- that it passes through every canon event, preserves the required facts, and leaves each option's keyed profile unchanged (NFR-15) --- and one that fails is regenerated or replaced by an instance from the approved sample. The content administrator does not author occupational content. The administrator manages the occupation as an entry in the catalogue, arranges how approved content is laid out and presented in gameplay, reviews the occupation as a whole before release, publishes it, and operates its lifecycle thereafter. Published content is revalidated annually by the experts, since occupational practice changes. In the present phase, in which no industry experts are yet engaged, the project team performs the expert role on provisional seeds (Section 4.2.1).

The evaluation process is internal and supports the research objectives. It proceeds in two stages (Section 4.2.2): a synthetic benchmark constructed by the project team, and a reference set of real responses and generated scenarios rated independently by domain experts, against which generation quality and assessment accuracy are measured offline as specified in Section 2.6.

\subsubsection{Business rules}

Business rules are stated here as policies and constraints of the domain that hold independently of any particular implementation. Constraints arising from implementation choices --- which models are used, how cost is bounded, where data is processed --- are recorded separately as design constraints in Section 4.3.3.

{\small
\begin{longtable}{|>{\raggedright\arraybackslash}p{\dimexpr 0.103\textwidth-2\tabcolsep-\arrayrulewidth\relax}|>{\raggedright\arraybackslash}p{\dimexpr 0.573\textwidth-2\tabcolsep-\arrayrulewidth\relax}|>{\raggedright\arraybackslash}p{\dimexpr 0.319\textwidth-2\tabcolsep-\arrayrulewidth\relax}|}
\caption{Business rules governing the proposed system.} \\
\hline
\textbf{ID} & \textbf{Rule} & \textbf{Origin} \\ \hline
\endfirsthead
\hline
\textbf{ID} & \textbf{Rule} & \textbf{Origin} \\ \hline
\endhead
BR-01 & The entry level of every occupation is accessible without prerequisite. & Exploration is the purpose; gating discovery defeats it \\ \hline
BR-02 & Access to a higher level requires the competence credit specified for that level. & Competency-based progression (Section 2.2) \\ \hline
BR-03 & Demonstrated competencies that overlap an adjacent occupation's requirements, as determined from the occupational taxonomy, count towards that occupation's levels, so that a user may enter it above its entry level. & Sections 2.2, 2.6; a user exploring related work need not restart \\ \hline
BR-04 & An occupation may be rated only by a user who has completed it. & Rating validity depends on exposure \\ \hline
BR-05 & No scenario seed or rubric is published without domain-expert validation of the seed and of a sample of the scenarios generated from it. & Section 2.5; accuracy of occupational representation \\ \hline
BR-06 & Assessment results are formative. They are not transmitted to employers and do not determine access to any external opportunity. & Section 2.7; proportionality to assessment reliability \\ \hline
BR-07 & Occupational representation must include the unfavourable aspects of the work. & Realistic job preview evidence (Section 2.1) \\ \hline
BR-11 & Replaying a scenario cannot raise a competency beyond the credit that scenario can award. & Keeps unlocks tied to demonstrated competence (BR-02) \\ \hline
BR-13 & An occupation's content becomes available to explorers only when the content administrator publishes it, which requires approved seeds (BR-05). & Scenario lifecycle (Section 5.1.1) \\ \hline
BR-14 & Revising a seed withdraws the scenarios generated from it until the revised seed is approved. & Scenario lifecycle; accuracy of occupational representation \\ \hline
BR-15 & Published content is revalidated by domain experts every year. & Occupational practice changes \\ \hline
BR-16 & A retired occupation remains in the records of users who attempted it but can no longer be started. & Scenario lifecycle; continuity of users' records \\ \hline
\end{longtable}
}

BR-01 and BR-03 are complementary. BR-01 keeps discovery free: the entry level of any occupation can be tried without prerequisite. BR-03 concerns what earlier work is worth: competence already demonstrated in one occupation is credited towards an adjacent one, so that the user does not begin it again from the entry level.

\subsubsection{Organisational policies and legal-ethical considerations}

The platform processes personal data, including performance data, of users who may be minors, and is therefore subject to Vietnam's Law on Personal Data Protection and its implementing decree\cite{r99,r100} and to comparable frameworks for users in other jurisdictions. Three policies follow. Data minimisation limits collection to what the exploration and assessment functions require. Aggregate-only publication governs any statistic derived from user performance, with suppression thresholds applied so that small populations cannot be de-anonymised; as a privacy measure rather than a rule of the domain, it is specified as a non-functional requirement (NFR-06). Formative-use limitation (BR-06) prevents assessment output from acquiring consequences the method's reliability does not support.

Two provisions of the law shape the design directly. A child's rights as a data subject are exercised by the child's legal representative; the platform therefore obtains guardian consent for users under sixteen, as a conservative practice. Sending a user's response to a model provider located abroad is a cross-border transfer of personal data, for which the law requires a transfer impact assessment to be filed with the competent authority within sixty days of the first transfer; responses are therefore stripped of identifying details before they leave the platform, since de-identified data is not personal data under the law, and the assessment is to be prepared before any deployment to real users. How these provisions apply to a university-hosted project, and whether behavioural logs fall within the decree's sensitive categories, are to be confirmed with the university's legal office.

Two further policies concern the system's use of generative models. Provenance disclosure requires that users be informed that scenarios are machine-generated from expert-authored material and that assessments are machine-produced. Expert accountability assigns responsibility for the occupational accuracy of published content to the authoring and reviewing domain experts, as recorded in the system, rather than to the generating model. Under the intended partnership described in Section 4.1.5, data governance would additionally follow the requirements of the Ministry of Education and Training for systems serving students. Detailed processing policies, including third-party model providers, are specified in Section 5.3.10.

\subsubsection{Business model and organisation}

JobQuest is designed as a public-benefit platform rather than a commercial product, and its business model follows from that position and from the evidence on which the system rests.

Value and beneficiaries. The platform offers learners experiential occupational information before commitment, and offers education authorities and institutions a scalable complement to limited counselling capacity. All content is available to all users without charge; access to levels and occupations is governed by demonstrated competence (BR-01 to BR-03) rather than by payment. The choice is deliberate: the users who most need occupational information --- secondary students choosing a field under examination and family pressure --- are those least able to pay for it, and the platform's effect depends on reaching them.

Partners and funders. The Ministry of Education and Training and educational-technology organisations may each act both as partners and as funders. As a partner, the Ministry --- preceded by city-level education departments --- supplies distribution through schools, legitimacy with schools and families, and a governance framework for minors' data, while edtech organisations supply pilot hosting, distribution, and technical cooperation. As funders, both may provide programme funding, grants, or in-kind support such as cloud credits. Industry partners and professional associations supply and coordinate domain experts, and the host university provides academic oversight. One set of safeguards applies to every partner and funder, whatever its role: none authors, validates, or ranks occupational content; none receives identifiable user data, only aggregate statistics (NFR-06); none may introduce charges, promotional content, or monetisation of users; and every partnership and source of funding is disclosed on the platform.

Funding and sustainability. Without user revenue, operation is funded through public and philanthropic support: grant funding sought from the Ministry and from educational-technology organisations; sponsorship from companies, which may contribute experts as well as funds; voluntary donations through a donation box on the platform; and in-kind contribution of expert consultation. The first expert panel, in Stage 1, contributes voluntarily; from Stage 2 an honorarium for experts is budgeted within grant funding. The cost structure determines how that funding must be managed. The dominant variable cost is language-model inference, incurred per session, which is why per-session cost is a monitored requirement (NFR-11). The system manages it by enforcing a configured ceiling on the model cost of each session, a cost-management mechanism specified with the design in Chapter 5. The dominant fixed cost is expert authoring, incurred per occupation, which is why coverage expands field by field rather than all at once. Model access is provider-portable (NFR-12, DC-02), so that inference can move to lower-cost or institutionally hosted models as funding conditions change. Two further levers reduce cost. A panel of smaller judging models from different families is less expensive than a single large judge\cite{r74}, and because a seed fixes a scenario's canon events, keyed options, and rubric in advance (BR-05), run-time generation is confined to the narrative between those events, which bounds the length and therefore the cost of each generation.

Sponsorship and content integrity. A model in which employers paid to shape how occupations are portrayed would reward favourable portrayal, in direct conflict with the realistic job preview evidence (Section 2.1) and with BR-07. Sponsorship is therefore admitted only as acknowledgement. Sponsoring companies are named, with their logos, under a plain ``Supported by'' heading; the acknowledgement carries no slogans, product claims, or links to product pages, and it never appears within scenarios, assessments, or rankings. Sponsors do not author, approve, or rank content, and every sponsorship is disclosed. Whether such an acknowledgement falls within the advertising regulations is subject to legal review, to be settled with the university's legal or partnership office before any logo is displayed.

Perception of a free service. A service offered without charge can be taken to be of low quality. The platform answers this with signals of credibility rather than with a price. Sponsor acknowledgements and the donation box are present from the outset. A badge marks scenarios validated by domain experts; it is shown only on scenarios that experts have actually approved, and therefore from Stage 1. Evaluation results are published once Stage A or Stage B has produced them, with the stage stated, since Stage A does not establish agreement with experts. Endorsement by the Ministry or by partner institutions is deferred to Stage 3. No paid tier is introduced.

Starting point and growth roadmap. A public authority does not adopt an unproven platform, so the project begins where it stands and advances only as evidence accumulates. The policy direction supports the end goal: the national scheme on career guidance and student streaming in general education, approved by Decision 522/QĐ-TTg, covered the period 2018--2025\cite{r95}; its results were reviewed in November 2025, with a decree on career guidance and educational streaming in preparation\cite{r96}; and Ho Chi Minh City has set a 2030 target of at least 35\% of students in basic sciences, engineering, and technology\cite{r97}. Table 4.2 sets out the stages. Each has a gate --- the evidence required before the next begins --- and meeting the thresholds of the two research problems is the evidence that justifies approaching funders and authorities.

{\small
\begin{longtable}{|>{\raggedright\arraybackslash}p{\dimexpr 0.171\textwidth-2\tabcolsep-\arrayrulewidth\relax}|>{\raggedright\arraybackslash}p{\dimexpr 0.159\textwidth-2\tabcolsep-\arrayrulewidth\relax}|>{\raggedright\arraybackslash}p{\dimexpr 0.301\textwidth-2\tabcolsep-\arrayrulewidth\relax}|>{\raggedright\arraybackslash}p{\dimexpr 0.193\textwidth-2\tabcolsep-\arrayrulewidth\relax}|>{\raggedright\arraybackslash}p{\dimexpr 0.171\textwidth-2\tabcolsep-\arrayrulewidth\relax}|}
\caption{Staged roadmap from university project to national public resource.} \\
\hline
\textbf{Stage} & \textbf{Indicative period} & \textbf{Activities} & \textbf{Resources and funding} & \textbf{Gate to next stage} \\ \hline
\endfirsthead
\hline
\textbf{Stage} & \textbf{Indicative period} & \textbf{Activities} & \textbf{Resources and funding} & \textbf{Gate to next stage} \\ \hline
\endhead
0. University project & Present phase, to the end of 2026 & Core system for the IT field; provisional seeds prepared by the team from authoritative public sources under advisor guidance; Stage A evaluation & Project team; advisor; host university; model costs self-funded or covered by cloud credits & Working system; Stage A targets met \\ \hline
1. Capstone and first pilot & Capstone phase, 2027 & Complete product; first expert panel for three pilot roles (back-end engineering, DevOps and cloud engineering, and business analysis); pilot with volunteer IT students; Stage B evaluation & University research or innovation support; cloud credits; competitions; donations & RP-1 and RP-2 thresholds met for the pilot roles \\ \hline
2. Incubation and multi-site pilot & 2027--2028 & Operating organisation established; expert panels extended through industry partners; pilots with partner schools and university career centres in Ho Chi Minh City & Edtech grants; company sponsorship; donations; in-kind expert consultation & Multi-site evidence of use and decision clarity; expert coverage of the IT roles \\ \hline
3. Public-sector partnership & 2028--2029 & Evidence presented to the city Department of Education and Training, then to the Ministry; alignment with the successor to Scheme 522; a second occupational field begins & Public programme funding; edtech grants; company sponsorship; donations & A formal cooperation agreement \\ \hline
4. National public resource & 2030 onward & Deployment through the public education system; further fields; annual expert revalidation & Public and philanthropic funding; company sponsorship; donations & Sustained operation \\ \hline
\end{longtable}
}

The project remains university-hosted during Stages 0 and 1, which gives it institutional standing without a separate legal entity. The operating organisation formed in Stage 2 is to take a non-commercial form consistent with free access, such as a social enterprise or a university-affiliated centre, to be settled with legal advice at that stage.

Organisational structure. The organisation is staged. During the present phase and the capstone phase, the three-member project team holds every operating role, one member holding several where necessary; domain experts join from the capstone phase; and the project advisor provides academic oversight. For sustained operation after the capstone phase, the proposed structure places a steering committee --- with representatives of the Ministry, the host university, funding organisations, and the project lead --- above five operating units, advised by an expert board of senior domain and assessment specialists. Table 4.3 sets out the units, their positions, and their responsibilities. The indicative core is about eleven full-time positions, with domain leads part-time and contributing experts engaged per occupation, a size chosen to remain sustainable on grant funding. The numbers of positions are indicative and the structure is not rigid: in the early stages one person may hold more than one position.

{\small
\begin{longtable}{|>{\raggedright\arraybackslash}p{\dimexpr 0.182\textwidth-2\tabcolsep-\arrayrulewidth\relax}|>{\raggedright\arraybackslash}p{\dimexpr 0.271\textwidth-2\tabcolsep-\arrayrulewidth\relax}|>{\raggedright\arraybackslash}p{\dimexpr 0.338\textwidth-2\tabcolsep-\arrayrulewidth\relax}|>{\raggedright\arraybackslash}p{\dimexpr 0.204\textwidth-2\tabcolsep-\arrayrulewidth\relax}|}
\caption{Proposed organisational structure for sustained operation.} \\
\hline
\textbf{Unit} & \textbf{Positions (indicative number)} & \textbf{Responsibilities} & \textbf{Corresponding user group or constraint} \\ \hline
\endfirsthead
\hline
\textbf{Unit} & \textbf{Positions (indicative number)} & \textbf{Responsibilities} & \textbf{Corresponding user group or constraint} \\ \hline
\endhead
Product and Engineering & Engineering lead (1); full-stack engineers (2); AI engineer (1); DevOps and system administrator (1) & Platform development and operation; AI model integration, configuration, and cost control; infrastructure and security & System administrator; DC-01, DC-02, DC-05; NFR-11 \\ \hline
Content and Occupational Expertise & Content lead (1); assessment and rubric specialist (1); domain leads, one per occupational domain (part-time); domain experts (part-time, three to four per role) & Occupation catalogue; authoring of seeds, rubrics, and keyed options by domain experts, and their review of generated scenarios; presentation, review, and publication of occupations; annual revalidation & Domain expert; content administrator; BR-05, BR-07 \\ \hline
Research and Evaluation & Research lead (1); evaluation analyst (1) & Construction of the expert reference set; offline evaluation; subsequent user studies & Evaluation process; RP-1, RP-2 \\ \hline
Partnerships and Outreach & Partnerships manager (1); school and institution outreach coordinator (1) & Liaison with the Ministry, schools, and universities; grant applications and reporting; recruitment of domain experts & --- \\ \hline
Data Protection and Compliance & Data protection officer (1, part-time initially) & Consent and minors' data; data-subject requests; regulatory compliance & DC-03; policies in Section 4.1.4 \\ \hline
\end{longtable}
}

Occupational domains are the unit of content governance: each has a domain lead accountable for its accuracy. Within the information technology field covered in the present phase, the proposed domains are software engineering; data and artificial intelligence; network and systems infrastructure; information security; and IT product and business analysis. Each further field is added as a new domain with its own lead, which is how coverage expands without restructuring the organisation.

\subsection{System Description}

\subsubsection{System overview}

JobQuest is a web-based platform that presents occupations as sequences of performable tasks, assesses the competencies a user demonstrates while performing them, and uses those assessments to structure further exploration. Its purpose is to supply occupational information experientially, at a breadth that authored content cannot reach, to people who must choose or change an occupation.

The experience is delivered through three presentation and progression conventions selected against the evidence in Section 2.3 rather than by preference. Progression follows a role-playing convention, in which levels and access to further occupations are earned through demonstrated competence rather than time spent, which the role-play evidence supports as the genre most associated with skill acquisition. Scenarios are presented in a branching narrative form comparable to a visual novel, the narrative moderator being among the strongest identified in the gamification meta-analyses. Career and life events are introduced probabilistically across sessions, a convention drawn from roguelike design, so that the circumstances bearing on occupational fit vary between runs; consistent with the harm evidence in Section 2.3, this randomisation governs narrative variety only and is not coupled to any reward schedule or monetisation. These conventions carry engagement; consistent with Section 2.8 they are not relied upon for competence development, which rests on task authenticity and assessment.

The system serves four user groups: the explorer, the domain expert, the content administrator, and the system administrator.

The explorer is the primary user: a student choosing a field of study, a graduate entering the labour market, or a worker considering transition. The explorer performs tasks, receives assessment and guidance, accumulates competence credit, and navigates among occupations. All other groups exist to make this possible.

The content administrator manages the occupation catalogue and is responsible for how expert-authored content reaches explorers: arranging its layout and presentation in gameplay; creating, reviewing, and publishing occupations; and operating the content lifecycle, including monitoring coverage and currency across occupations and arranging revalidation. The content administrator does not author or alter occupational content, which belongs to the domain experts.

The system administrator operates the platform: configuring and monitoring the generation and assessment models, managing capacity and cost, administering accounts and access, and overseeing data protection controls.

Domain experts are senior or lead practitioners in an occupation, typically with around five or more years in the role and experience supervising junior staff, who therefore know what entry-level work demands and where newcomers typically fail. They use the system directly: they author and own the occupational content of scenarios --- the required facts, target competencies, response options and their effectiveness keys, and rubrics --- and they review the scenarios the system generates and the cases in which automated judgment departs from their own. Each role is served by three to four experts, three being the accepted minimum for content validation\cite{r69}, drawn from at least two types of organisation --- such as product, outsourcing, startup, and in-house enterprise IT --- because the same job title involves different work in different organisations; no expert reviews scenarios generated from a seed they helped prepare. Their involvement is heaviest when an occupation's seeds are first prepared and lighter thereafter, through annual revalidation, and is arranged part-time through partner organisations; the first panel is recruited through the advisor's and the university's alumni networks and contributes voluntarily. Across the approximately twenty-two IT roles in scope, the full network is some sixty-six to eighty-eight experts, reached progressively (Section 4.1.5). In the present phase the project team prepares provisional seeds from authoritative public sources under the advisor's guidance, and these are not presented as expert-validated. This group is the guarantor of the system's central claim, that its representations correspond to real work.

\subsubsection{Research problems}

In this report, a research problem is a question that meets three conditions: the answer is not established in the existing literature; the answer constitutes knowledge that generalises beyond this particular system; and the project's evaluation produces evidence that answers it. A question resolved by a design decision fails the second condition, and a question about how evidence is gathered concerns method rather than a problem. Two questions meet all three conditions, both in artificial intelligence applied to education, and they correspond to the first two specific objectives in Section 1.2.

RP-1, grounded generation of occupational scenarios. The literature establishes that language models generate fluent instructional content but that fidelity to a specified domain requires explicit constraint (Section 2.5); it does not establish how accurate expert-seeded generation is for task-based occupational content, nor whether generated tasks preserve the competency they target (Section 2.2). The question is therefore: when generation is constrained by an expert-authored seed specifying required facts and target competencies, what proportion of generated scenarios do domain experts judge to be occupationally accurate and to require the target competency, and how does this compare with generation from the same instructions without the seed? The comparison is what makes the answer knowledge rather than description, since without it an accuracy figure cannot be attributed to the seeding.

RP-2, validity of language-model stealth assessment. Reported agreement between automated judges and human raters is systematically optimistic, and chance-corrected, difficulty-stratified validation is rarely reported (Section 2.5); nor has such assessment been validated where the evidence of competence is the performance of a task rather than an answer to a test item (Section 2.2). The question is therefore: to what extent does rubric-based automated judgment of task responses agree with domain-expert judgment of the same responses, measured by quadratic weighted kappa, relative to the agreement between experts themselves, and how does that agreement vary with item difficulty and across repeated runs? Expert--expert agreement provides the ceiling against which the automated judge is interpreted.

Both problems are evaluated in two stages, because the experts on whom they depend are engaged progressively (Section 4.2.1). In Stage A, during the present phase, the project team constructs a synthetic benchmark: for each free-text task, a strong reference response and variants that each degrade one rubric criterion to a known level, together with probes for length, format, and ordering bias, following the perturbation-checklist method\cite{r73}. A panel of judging models from different families, none of which produced the synthetic responses, scores them independently and repeatedly\cite{r74}, and the advisor confirms the designed levels on a random sample of about a fifth of the variants. Stage A establishes whether the judge is sensitive to each criterion, robust to known biases, and consistent across runs; it cannot establish agreement with experts on real responses, and is not presented as doing so. In Stage B, on the completed product, volunteer participants complete tasks with informed consent, and their responses, together with generated scenarios, are rated independently and blind by three experts per pilot role. RP-1 is then answered by content validity indices, with I-CVI of at least .78 and S-CVI/Ave of at least .90\cite{r70}, by a paired comparison of seeded and unseeded generation using McNemar's test on acceptability, and by a fact-preservation target of 95\% of seed-required facts. RP-2 is answered against the three acceptance criteria of Williamson, Xi and Breyer\cite{r71}, with agreement among experts measured by Krippendorff's alpha\cite{r72}.

Three further questions arise in the project and are addressed, but they are not research problems under these conditions.

{\small
\begin{longtable}{|>{\raggedright\arraybackslash}p{\dimexpr 0.358\textwidth-2\tabcolsep-\arrayrulewidth\relax}|>{\raggedright\arraybackslash}p{\dimexpr 0.471\textwidth-2\tabcolsep-\arrayrulewidth\relax}|>{\raggedright\arraybackslash}p{\dimexpr 0.167\textwidth-2\tabcolsep-\arrayrulewidth\relax}|}
\caption{Questions addressed in the project but not treated as research problems.} \\
\hline
\textbf{Question} & \textbf{Condition not met} & \textbf{Where it is addressed} \\ \hline
\endfirsthead
\hline
\textbf{Question} & \textbf{Condition not met} & \textbf{Where it is addressed} \\ \hline
\endhead
How should demonstrated competence determine access to adjacent occupations? & Resolved by a design decision based on occupational taxonomies; the project does not generate new knowledge about it & Design (Chapter 5) \\ \hline
Does the engagement design sustain exploration without displacing learning? & Answerable only through longitudinal user study, which the offline evaluation does not provide & Subsequent work \\ \hline
How should the expert reference set be constructed? & Concerns how evidence is gathered --- a matter of method rather than a problem & Evaluation (Chapter 7) \\ \hline
\end{longtable}
}

\subsubsection{Challenges and difficulties}

Generation fidelity is the principal technical difficulty. Occupational content is factual, and a scenario that is plausible but inaccurate is worse than no scenario, since it conveys a confident false impression of a real profession. Long-horizon consistency compounds this within multi-task levels.

Assessment validity is the principal methodological difficulty. Section 2.5 indicates that agreement figures for automated judgment are frequently overstated and that performance degrades unevenly with difficulty; obtaining defensible figures requires a reference set that is stratified and expert-authored, which is expensive to produce.

Expert availability constrains breadth. Each role requires three to four senior practitioners for authoring and review, whose time must be arranged through partner organisations; in Stage 1, an estimated seven hours per expert covers scenario ratings, response ratings, and option keying for a pilot role. The three pilot roles are back-end engineering, DevOps and cloud engineering, and business analysis, the last chosen so that assessment is also examined on work whose output is requirements and communication with stakeholders rather than code. Because experts are engaged progressively, the present phase relies on provisional seeds prepared from public sources, and occupational accuracy can be claimed only once the Stage 1 panel has reviewed them. This is the reason for restricting the first phase to a single field, and for treating expert authoring as a first-class activity of the system rather than an administrative convenience.

Cold start affects the aggregate occupational profiles, which require a population of completions before they are informative. Until that population exists, profiles are absent or unreliable, and the interface must represent this honestly rather than display an unfounded figure.

Evaluation without deployment limits what the project can establish. Offline benchmarking substantiates generation quality and assessment accuracy but cannot establish engagement retention or effects on career decisions, which require longitudinal study. The claims are bounded accordingly, consistent with Section 2.8.

Operating without user revenue makes cost a design constraint rather than a commercial concern. Inference cost per session must remain within what grant funding can sustain, which constrains model choice, response length, and caching strategy.

Interaction modality limits what the platform can represent, and the limit differs by competency. Hamstra et al.\cite{r98} argue that the realism of a simulation should be understood not as physical resemblance but as functional task alignment --- the correspondence between what a task demands and what the real situation demands --- noting that physical fidelity has been found largely independent of educational effectiveness. On this view the relevant question is not whether a screen-based task looks like a workplace but whether it demands what the work demands. Table 4.5 compares the competencies a JobQuest task can involve with their counterparts in a live work situation, rating alignment on four levels: high, moderate, partial, and low.

{\small
\begin{longtable}{|>{\raggedright\arraybackslash}p{\dimexpr 0.267\textwidth-2\tabcolsep-\arrayrulewidth\relax}|>{\raggedright\arraybackslash}p{\dimexpr 0.338\textwidth-2\tabcolsep-\arrayrulewidth\relax}|>{\raggedright\arraybackslash}p{\dimexpr 0.231\textwidth-2\tabcolsep-\arrayrulewidth\relax}|>{\raggedright\arraybackslash}p{\dimexpr 0.159\textwidth-2\tabcolsep-\arrayrulewidth\relax}|}
\caption{Alignment of competencies between screen-based tasks and live work situations.} \\
\hline
\textbf{Competency} & \textbf{In a JobQuest task} & \textbf{In a live work situation} & \textbf{Alignment} \\ \hline
\endfirsthead
\hline
\textbf{Competency} & \textbf{In a JobQuest task} & \textbf{In a live work situation} & \textbf{Alignment} \\ \hline
\endhead
Technical analysis and problem-solving, such as diagnosing a failing service & Performed through the same text and code interfaces used in practice & Largely screen-based in IT roles & High \\ \hline
Written communication, such as incident reports, documentation, and messages to colleagues & Produced and assessed as text & Produced as text & High \\ \hline
Prioritisation and decision-making under constraints & Decisions made with consequences represented in the scenario & The same decisions, with real consequences for users and colleagues & Moderate: decision structure aligned, stakes attenuated \\ \hline
Adaptability to changing requirements & Changes introduced as scenario events & Changes arrive unpredictably and compound & Moderate \\ \hline
Collaboration and handling of stakeholders & Text interaction with AI-driven colleagues and clients & Real people with their own goals, relationships, and history & Partial: structure aligned, social dynamics simplified \\ \hline
Working under time pressure and regulating stress & Time limits can be simulated, with low personal stakes & Sustained pressure with genuine stakes & Low \\ \hline
Oral communication and presentation & Represented only through what is written; delivery, tone, and listening are not exercised & Central in many roles & Low \\ \hline
Interviews & The content of answers can be rehearsed in text; live delivery under observation is not exercised & The gateway to most roles & Low \\ \hline
Physical and manual skill & Procedures can be described and sequenced; manual execution is not exercised & Central in some occupations & Low \\ \hline
\end{longtable}
}

Two consequences follow. Alignment is highest for the competencies at the core of information technology work, which is itself largely screen-mediated --- a further reason for the first phase's coverage. Conversely, competencies that depend on live interpersonal dynamics, genuine stakes, speech, or physical performance are attenuated: oral communication, interviewing, and physical skill are represented only through their written or procedural content, and alignment for them is low. Both the assessment and the self-efficacy the platform builds (Section 2.4) are therefore evidence about the well-aligned competencies only, and results must not be presented as indicating readiness for the others. Voice interaction could extend coverage to oral communication; genuine stakes and physical skill lie outside what any screen-based simulation can provide.

Localisation presents a further difficulty: occupational taxonomies are authored in English for other labour markets, and their occupational structures do not map cleanly onto Vietnamese practice, requiring expert adaptation rather than translation.

\subsection{User's Requirements}

Requirements are expressed first as user requirements per group, each with a priority (Must, Should, Could). These are then analysed into functional, non-functional, and data requirements. Each of these carries two verification methods: a business verification, by which a stakeholder --- an explorer, a domain expert, a content administrator, or a system administrator --- confirms that the requirement meets its purpose, through acceptance tests, walkthroughs, or expert review; and a technical verification, by which the development team confirms that the implementation behaves correctly, through automated, integration, boundary, load, or fault-injection tests.

\subsubsection{User requirements}

User requirements --- explorer:

{\small
\begin{longtable}{|>{\raggedright\arraybackslash}p{\dimexpr 0.114\textwidth-2\tabcolsep-\arrayrulewidth\relax}|>{\raggedright\arraybackslash}p{\dimexpr 0.744\textwidth-2\tabcolsep-\arrayrulewidth\relax}|>{\raggedright\arraybackslash}p{\dimexpr 0.137\textwidth-2\tabcolsep-\arrayrulewidth\relax}|}
\caption{User requirements of explorers.} \\
\hline
\textbf{ID} & \textbf{User requirement} & \textbf{Priority} \\ \hline
\endfirsthead
\hline
\textbf{ID} & \textbf{User requirement} & \textbf{Priority} \\ \hline
\endhead
UR-E01 & As an explorer, I want to create an account, recover it if I forget my password, and return to my own progress on every visit, so that my exploration is kept safe. & Must \\ \hline
UR-E02 & As an explorer, I want to browse the occupations available, read what each involves at every level, including its unfavourable aspects, and try the entry level of any of them without prerequisite, so that I can survey possibilities before committing to one. & Must \\ \hline
UR-E03 & As an explorer who does not know where to begin, I want optional suggestions of occupations to start with, and a description of myself that explains them, so that a blank choice does not stop me. & Should \\ \hline
UR-E04 & As an explorer, I want to be told my role, team, and working context before each scenario begins, so that my decisions are situated rather than abstract. & Must \\ \hline
UR-E05 & As an explorer, I want to perform tasks that practitioners would recognise and whose objective I can understand, so that what I learn reflects the occupation. & Must \\ \hline
UR-E06 & As an explorer, I want to respond in the forms the work itself takes --- choosing a course of action, ordering steps, prioritising tasks, or writing an answer such as an email or a code review --- so that I experience how the work is done, not only what it is about. & Must \\ \hline
UR-E07 & As an explorer, I want hints when I am stuck, at a small cost to my best possible score, so that I can continue while still learning. & Should \\ \hline
UR-E08 & As an explorer, I want my response judged as an expert in the occupation would judge it, probed further when my answer is unclear, and returned as a result, an explanation, and guidance, so that I can trust the judgment and know how to improve. & Must \\ \hline
UR-E09 & As an explorer, I want follow-up questions and feedback to arrive without long waits, and to see progress while I wait, so that the experience keeps its momentum. & Must \\ \hline
UR-E10 & As an explorer, I want to play on my phone, over a weak connection, and with assistive technology if I need it, without installing anything, so that I can explore wherever I am. & Must \\ \hline
UR-E11 & As an explorer, I want to resume where I stopped after an interruption or a system failure, and to be told promptly if something has failed, so that I never lose work I have done. & Must \\ \hline
UR-E12 & As an explorer, I want to know which content and judgments are produced by AI, so that I can weigh them accordingly. & Must \\ \hline
UR-E13 & As an explorer, I want to choose scenarios within a level, have the competencies I demonstrate recorded, and open higher levels as they grow, without replaying the same scenario being enough to inflate them, so that my progress reflects what I can do rather than time spent. & Must \\ \hline
UR-E14 & As an explorer, I want to start an adjacent occupation at a level my competencies support, so that I need not restart from the entry level when exploring related work. & Should \\ \hline
UR-E15 & As an explorer, I want to see the career path across the levels of an occupation --- the role held at each level --- so that I understand where the occupation could lead me. & Should \\ \hline
UR-E16 & As an explorer, I want to rate an occupation after completing it and see how others rated it, without anyone's individual rating being exposed, so that I can compare my impression with a broader one. & Should \\ \hline
UR-E17 & As an explorer, I want a private record of what I have attempted and demonstrated, so that my exploration accumulates into something I can reflect on. & Should \\ \hline
UR-E18 & As an explorer, I want my individual responses and results never to be shown to others or passed to employers, and my data to be deleted when I ask, so that trying an occupation carries no risk to me. & Must \\ \hline
UR-E19 & As an explorer, I want to use the platform in Vietnamese or English, so that language is not a barrier to exploring. & Should \\ \hline
UR-E20 & As an explorer, I want short side quests I can answer on the spot for bonus points, so that I can engage with an occupation in small steps. & Could \\ \hline
UR-E21 & As an explorer, I want to send feedback about the whole app, so that problems and ideas reach the team. & Should \\ \hline
\end{longtable}
}

User requirements --- domain expert:

{\small
\begin{longtable}{|>{\raggedright\arraybackslash}p{\dimexpr 0.114\textwidth-2\tabcolsep-\arrayrulewidth\relax}|>{\raggedright\arraybackslash}p{\dimexpr 0.744\textwidth-2\tabcolsep-\arrayrulewidth\relax}|>{\raggedright\arraybackslash}p{\dimexpr 0.137\textwidth-2\tabcolsep-\arrayrulewidth\relax}|}
\caption{User requirements of domain experts.} \\
\hline
\textbf{ID} & \textbf{User requirement} & \textbf{Priority} \\ \hline
\endfirsthead
\hline
\textbf{ID} & \textbf{User requirement} & \textbf{Priority} \\ \hline
\endhead
UR-D01 & As a domain expert, I want to write a scenario's context with AI assistance and set the allowed question types, the score of each option, the grading rubric, hints, random events, and side quests directly in the system, with my authorship recorded, so that the content I own is traceable to me and needs no developer to enter. & Must \\ \hline
UR-D02 & As a domain expert, I want every published scenario to keep the required facts and option scores exactly as I set them, so that the system never misrepresents the work in my name. & Must \\ \hline
UR-D03 & As a domain expert, I want to review scenarios written by another expert of my occupation against shared criteria and approve them, request changes, or reject them, so that only content experts judge accurate, and AI output kept within the limits they set, reaches explorers. & Must \\ \hline
UR-D04 & As a domain expert, I want to see where the automated judge disagrees with expert judgments, so that I can refine the rubric. & Should \\ \hline
UR-D05 & As a domain expert, I want to revalidate published content when practice changes, so that what explorers see stays current. & Should \\ \hline
UR-D06 & As a domain expert, I want to create and keep up to date the knowledge of an occupation once the content administrator has created it --- levels, salary, generation rules, competencies, and adjacency --- so that progression and adjacency reflect how the work is organised in Vietnam. & Should \\ \hline
UR-D07 & As a domain expert, I want to rate explorers' responses and generated scenarios blind, independently of other experts, so that the automated judge and the generator can be measured against expert judgment. & Should \\ \hline
UR-D08 & As the author of a scenario, I want to play the draft game before it is published, without my play being counted as explorer data, so that I can confirm it presents my content correctly. & Must \\ \hline
UR-D09 & As the author of a scenario, I want to add sample answers with expert scores while preparing it, so that the automated judge can be checked against my judgment. & Should \\ \hline
\end{longtable}
}

User requirements --- content administrator:

{\small
\begin{longtable}{|>{\raggedright\arraybackslash}p{\dimexpr 0.114\textwidth-2\tabcolsep-\arrayrulewidth\relax}|>{\raggedright\arraybackslash}p{\dimexpr 0.744\textwidth-2\tabcolsep-\arrayrulewidth\relax}|>{\raggedright\arraybackslash}p{\dimexpr 0.137\textwidth-2\tabcolsep-\arrayrulewidth\relax}|}
\caption{User requirements of content administrators.} \\
\hline
\textbf{ID} & \textbf{User requirement} & \textbf{Priority} \\ \hline
\endfirsthead
\hline
\textbf{ID} & \textbf{User requirement} & \textbf{Priority} \\ \hline
\endhead
UR-C01 & As a content administrator, I want to add a new occupation and track it through authoring, review, publication, and retirement without asking developers, so that catalogue growth is managed. & Must \\ \hline
UR-C02 & As a content administrator, I want to arrange how approved content is presented in gameplay --- question types, characters, and scenes --- without altering its substance, so that explorers receive it in a coherent form. & Must \\ \hline
UR-C03 & As a content administrator, I want to publish a scenario only after it has been approved and play-tested by its author, so that nothing unvalidated is released. & Must \\ \hline
UR-C04 & As a content administrator, I want to see coverage and play statistics for my occupations, so that I can identify outdated or weak content and arrange its revalidation. & Should \\ \hline
UR-C05 & As a content administrator, I want to replace an outdated scenario by removing the old one first and then opening a new one, so that explorers never meet superseded content. & Should \\ \hline
UR-C06 & As a content administrator, I want to open a scenario for each task and assign a domain expert other than its author to review it, so that every scenario is reviewed independently. & Must \\ \hline
UR-C07 & As a content administrator, I want to import occupational reference data from O*NET and ESCO, so that occupations start from an established source. & Should \\ \hline
UR-C08 & As a content administrator, I want to define the criteria domain experts use to review scenarios, so that every review follows the same standard. & Must \\ \hline
\end{longtable}
}

User requirements --- system administrator:

{\small
\begin{longtable}{|>{\raggedright\arraybackslash}p{\dimexpr 0.114\textwidth-2\tabcolsep-\arrayrulewidth\relax}|>{\raggedright\arraybackslash}p{\dimexpr 0.744\textwidth-2\tabcolsep-\arrayrulewidth\relax}|>{\raggedright\arraybackslash}p{\dimexpr 0.137\textwidth-2\tabcolsep-\arrayrulewidth\relax}|}
\caption{User requirements of system administrators.} \\
\hline
\textbf{ID} & \textbf{User requirement} & \textbf{Priority} \\ \hline
\endfirsthead
\hline
\textbf{ID} & \textbf{User requirement} & \textbf{Priority} \\ \hline
\endhead
UR-S01 & As a system administrator, I want to configure which models generate and which assess, keep the two separate, and change provider without touching any authored content, so that assessment stays independent and the platform is not locked to one vendor. & Must \\ \hline
UR-S02 & As a system administrator, I want to monitor usage, cost, and failures, re-run failed AI jobs, and keep the cost of each session within budget, so that the service remains operable on grant funding. & Must \\ \hline
UR-S03 & As a system administrator, I want to manage every account, explorers included --- grant staff their roles and occupations, view activity logs, and disable or delete accounts --- so that each person can do exactly their part and misuse can be stopped. & Must \\ \hline
UR-S04 & As a system administrator, I want to handle data-subject requests and oversee data-protection controls, so that the platform meets its legal obligations to users, including minors. & Must \\ \hline
UR-S05 & As a system administrator, I want to review feedback sent about the app, so that recurring problems are found and fixed. & Should \\ \hline
\end{longtable}
}

\subsubsection{Functional requirements}

{\small
\begin{longtable}{|>{\raggedright\arraybackslash}p{\dimexpr 0.103\textwidth-2\tabcolsep-\arrayrulewidth\relax}|>{\raggedright\arraybackslash}p{\dimexpr 0.193\textwidth-2\tabcolsep-\arrayrulewidth\relax}|>{\raggedright\arraybackslash}p{\dimexpr 0.193\textwidth-2\tabcolsep-\arrayrulewidth\relax}|>{\raggedright\arraybackslash}p{\dimexpr 0.197\textwidth-2\tabcolsep-\arrayrulewidth\relax}|>{\raggedright\arraybackslash}p{\dimexpr 0.309\textwidth-2\tabcolsep-\arrayrulewidth\relax}|}
\caption{Functional requirements with business and technical verification.} \\
\hline
\textbf{ID} & \textbf{Functional requirement} & \textbf{Group} & \textbf{Business verification} & \textbf{Technical verification} \\ \hline
\endfirsthead
\hline
\textbf{ID} & \textbf{Functional requirement} & \textbf{Group} & \textbf{Business verification} & \textbf{Technical verification} \\ \hline
\endhead
FR-01 & Register, sign in, and recover the account & Explorer & Acceptance test: a user registers, signs in, resets a forgotten password, and reaches their own progress & Authentication tests; unauthenticated requests for progress are refused; reset links expire after one use \\ \hline
FR-02 & Edit, sign out of, and delete the account & Explorer & Walkthrough of each account action & Account deletion removes personal data as defined in NFR-05 \\ \hline
FR-03 & Run an optional, retakable RIASEC orientation quiz suggesting occupations\cite{r41} & Explorer & Pilot users complete the quiz and find the suggested occupations meaningful & Quiz yields a RIASEC profile; the five suggested occupations (project-defined) come from matching it with each occupation's O*NET interest code, without AI; skipping leads to the Galaxy Map; retaking replaces the earlier result \\ \hline
FR-04 & Write the explorer's portrait from the RIASEC profile with AI, including the social index & Explorer & Pilot users judge the description relevant and not misleading & The AI writes only the description text; the social index is the profile's Social score; the portrait is labelled as self-reported and not validated \\ \hline
FR-05 & Browse occupations on the Galaxy Map and show each profile with its downsides and level descriptions & Explorer & Explorers find an occupation on the map and can state its main downsides and what each level involves & Every published occupation appears on the map with all profile fields and a description for every level; ratings follow FR-19 suppression \\ \hline
FR-06 & Open the entry level of every published occupation to every explorer & Explorer & Acceptance test: a new explorer opens the entry level of any published occupation & Entry level reachable with zero credit and without any adjacency check \\ \hline
FR-07 & Let explorers start an adjacent occupation at a higher level when overlap conditions are met & Explorer & Domain experts judge the offered starting levels as plausible for seeded profiles & Overlap only unlocks higher starting levels and never blocks an entry level; no offer is made where adjacency data is missing \\ \hline
FR-08 & Enrol in levels, choose scenarios, award credit, and unlock the next level & Explorer & Acceptance test: an explorer enrols, chooses scenarios, and progresses through levels as expected & Threshold boundary tests unlock at and only at the specified credit; below the threshold, tasks for the missing competencies are suggested \\ \hline
FR-09 & Show the role, team, and context before each scenario & Explorer & Explorer walkthrough: the context is shown and understood before each scenario & End-to-end test: the context screen precedes every scenario; an unavailable scenario returns the explorer to the map \\ \hline
FR-10 & Support choice, ordering, prioritising, and free-text answers & Explorer & Explorers complete each response form without help & Known-input tests return the expected scores for each closed form \\ \hline
FR-11 & Score free-text answers with anchor points backed by quotations & Explorer & Domain experts judge a sample of scored responses as fair and well-evidenced & Every quotation is found verbatim in the response; scores outside the rubric anchors are rejected; agreement with experts is checked under NFR-14 \\ \hline
FR-12 & Answer follow-up questions generated by AI on free-text answers & Explorer & Explorers judge the follow-up questions relevant to what they wrote & Follow-ups stay within the expert's follow-up goal and the activity's limit; each follow-up passes the output check before display, otherwise the expert's default follow-up is shown; each follow-up answer is assessed \\ \hline
FR-13 & Offer hints when the explorer is stuck & Explorer & Explorers continue after a hint and still find the task meaningful & Using a hint removes the +2 option for that activity; hint use is recorded and shown in the summary; hints do not reveal grading criteria \\ \hline
FR-14 & Return result, feedback, and guidance separately & Explorer & Usability test: explorers identify the result, the feedback, and the guidance separately & All three present and distinguishable for every assessed scenario; AI-written feedback is built from the rubric anchors and passes the output check (on topic, no grading criteria, correct language), otherwise the expert's default feedback is shown; explorers can flag wrong content to the domain expert \\ \hline
FR-15 & Record evidenced expertise and soft-skill competencies, capped per skill per scenario & Explorer & Domain experts spot-check sampled sessions: recorded competencies match the observed behaviour & Records correspond to rubric criteria and are split into expertise and soft skills; random-event answers add bonus soft-skill evidence only; replay test confirms accumulation never exceeds the cap \\ \hline
FR-16 & Offer optional side quests for bonus points & Explorer & Explorers find side quests short and optional & Side quests are authored by the domain expert (FR-22); side-quest points are recorded separately from main-quest competence credit \\ \hline
FR-17 & Show the career path of each occupation level & Explorer & Explorers can describe the role held at each level of an occupation & Career path lists every level of the occupation with its role and description, as written by the domain expert \\ \hline
FR-18 & Show a private record of attempts, expertise and soft-skill points, and awards & Explorer & Explorers confirm their record is accurate and complete & Record reconstructs from session history for a seeded account; both point types shown in the top bar, summary, and profile; not accessible to other users \\ \hline
FR-19 & Collect occupation ratings and feedback & Explorer & Acceptance test: the rating and feedback form appears only after full completion; aggregated results are shown or withheld as expected & Form accepted only once every level is completed; aggregated result updates after each new submission; suppressed below the response-count threshold \\ \hline
FR-20 & Send app feedback and review it & Explorer, System administrator & An explorer sends feedback and a system administrator finds it in the review list & Every submission is stored with time and app version; the system administrator can mark it as handled \\ \hline
FR-21 & Create and update occupation knowledge: levels, salary, generation rules, competencies, adjacency & Domain expert & Experts approve the adaptation for Vietnamese practice & Knowledge can be created only for an occupation the content administrator has created; every rubric criterion resolves to a competency; routes without enough data are hidden, not guessed \\ \hline
FR-22 & Author scenario context, question types, option scores, rubric, hints, random events, and side quests & Domain expert & Walkthrough: an author prepares a scenario for an assigned task end to end, including its hints, random events, and side quests & Edits persist with authorship recorded; access limited to domain experts assigned to that occupation; every free-text activity has a follow-up goal, a default follow-up, and default feedback; random events are dropped into sessions at random within the configured probability, target the scenario's reviewed skills, and never lower the main score \\ \hline
FR-23 & Generate scenario content with AI from the author's context, keeping its generation history & Domain expert & Domain-expert review of sampled generations against the required facts and the occupation's rules & Automated validator passes on every stored generation; a scenario that fails the checks cannot be submitted; each run's context, model, cost, and check result are viewable \\ \hline
FR-24 & Add reference answers with expert scores & Domain expert & Experts confirm the stored criteria reflect what they intended & Item is stored only when validation passes; failing items return the reasons \\ \hline
FR-25 & Review another expert's scenario against the review criteria & Domain expert & Walkthrough: an expert approves, rejects, and requests changes on sample scenarios using the criteria & Reviewer is never the author; every criterion is answered before a decision; unapproved scenarios are never served to explorers; a rejected scenario cannot be reopened \\ \hline
FR-26 & Play-test the draft game before publication & Domain expert & Walkthrough: the author is invited, play-tests a draft game, and reports findings & Author is notified when the draft game is built; play-test sessions are excluded from every explorer aggregate; display findings return the scenario to presentation, content findings return it to the author and to review \\ \hline
FR-27 & Revalidate a published scenario on request & Domain expert & Experts revalidate a sample of published scenarios on request & A domain expert or content administrator can send a published scenario back to the review queue; the scenario stays published and playable meanwhile; the result is recorded; outdated scenarios are flagged for replacement \\ \hline
FR-28 & Report disagreements between AI and expert grading & Domain expert & Experts confirm the disagreement list helps them refine grading criteria & Disagreement list matches recomputation over the reference set \\ \hline
FR-29 & Rate responses and scenarios blind & Domain expert & Experts confirm they cannot see AI scores or other experts' ratings while rating & Blind-rating records exist per expert; inter-expert agreement is computable \\ \hline
FR-30 & Create and manage own occupations & Content administrator & Walkthrough: a content administrator takes their occupation from draft to published to retired & State transitions follow the defined lifecycle; publication requires completed review; actions on other occupations are refused; a domain that still holds occupations cannot be deleted \\ \hline
FR-31 & Import O*NET and ESCO reference data & Content administrator & A content administrator imports a sample and adjusts it & Imported records resolve against source identifiers; unresolved records are listed; domain experts are notified to derive competencies once the import completes \\ \hline
FR-32 & Define the scenario review criteria & Content administrator & Domain experts confirm the criteria are clear enough to apply & Reviews cannot start until criteria exist for the occupation; each review records an answer per criterion \\ \hline
FR-33 & Open one active scenario per task & Content administrator & Walkthrough: a content administrator opens a scenario and the author is notified & A second active scenario for the same task is refused; retired and rejected scenarios do not count \\ \hline
FR-34 & Assign a domain expert as reviewer, with a deadline, and reassign overdue reviews & Content administrator & Walkthrough: a content administrator assigns a domain expert to review a scenario and reassigns an overdue review & Assigning the author is refused; the domain expert is notified; the deadline appears in the review queue; the reviewer is reminded before the deadline; once it passes, the content administrator is notified and can reassign the review \\ \hline
FR-35 & Design the scenario presentation with AI & Content administrator & Walkthrough: a content administrator gamifies an approved scenario & Question types outside the allowed set are refused; content text is read-only in this step \\ \hline
FR-36 & Maintain a bias-checked asset library; AI generates missing assets & Content administrator & Reviewers judge sampled assets, including AI-generated ones, free of stereotyping & Every asset carries its source and bias-check result; an AI-generated asset enters the library only after its bias check passes \\ \hline
FR-37 & Publish a scenario after review and play-test & Content administrator & Walkthrough: a content administrator publishes a ready scenario with an effective date and release note & Publish is disabled unless the status is ready \\ \hline
FR-38 & Retire or delete a published scenario & Content administrator & Walkthrough: a content administrator replaces an outdated scenario & Played scenarios are retired and keep their history; unplayed scenarios are deleted permanently \\ \hline
FR-39 & Monitor own occupations: statistics and coverage & Content administrator & A content administrator spots a weak or outdated scenario from the dashboard & Figures match stored records; ratings follow the suppression threshold; coverage matches stored metadata \\ \hline
FR-40 & Assign and reassign one content administrator and domain experts to each occupation & System administrator & Walkthrough: a system administrator assigns staff to an occupation and moves it to another content administrator & Every occupation carries exactly one content administrator and at least one domain expert once assigned; reassigning an occupation moves its pending work to the new content administrator \\ \hline
FR-41 & Manage all user accounts, roles, and role-based access & System administrator & Walkthrough: a system administrator views the activity log of an explorer and a staff member, then disables and deletes an account & Staff sign in on the shared login screen and see only the work for their assigned occupations; wrong-role access to admin screens is refused; disabled accounts cannot sign in; deleted accounts follow NFR-05; staff-authored records are kept \\ \hline
FR-42 & Configure separate generation and assessment models & System administrator & A system administrator changes a provider through the configuration screen & Configuration rejects an identical provider for both roles \\ \hline
FR-43 & Monitor usage, cost, and errors, and re-run failed AI jobs & System administrator & A system administrator uses the monitoring screen to spot cost or error spikes & Metrics consistent with generated load; alerts raised when a cost threshold is crossed; failed generation and grading runs listed with cause and can be re-run \\ \hline
FR-44 & Handle data requests and data-protection settings & System administrator & Data-protection walkthrough of a deletion request & Requests are completed and logged; settings take effect on the next computation; a transfer impact assessment is documented and kept current for each AI provider \\ \hline
\end{longtable}
}

\subsubsection{Non-functional requirements}

{\small
\begin{longtable}{|>{\raggedright\arraybackslash}p{\dimexpr 0.114\textwidth-2\tabcolsep-\arrayrulewidth\relax}|>{\raggedright\arraybackslash}p{\dimexpr 0.182\textwidth-2\tabcolsep-\arrayrulewidth\relax}|>{\raggedright\arraybackslash}p{\dimexpr 0.213\textwidth-2\tabcolsep-\arrayrulewidth\relax}|>{\raggedright\arraybackslash}p{\dimexpr 0.221\textwidth-2\tabcolsep-\arrayrulewidth\relax}|>{\raggedright\arraybackslash}p{\dimexpr 0.265\textwidth-2\tabcolsep-\arrayrulewidth\relax}|}
\caption{Non-functional requirements with business and technical verification.} \\
\hline
\textbf{ID} & \textbf{Non-functional requirement} & \textbf{Category} & \textbf{Business verification} & \textbf{Technical verification} \\ \hline
\endfirsthead
\hline
\textbf{ID} & \textbf{Non-functional requirement} & \textbf{Category} & \textbf{Business verification} & \textbf{Technical verification} \\ \hline
\endhead
NFR-01 & Follow-up generation within 5 s (95th percentile); progress shown after 2 s & Performance & Usability session: explorers find the wait for follow-up questions acceptable and see progress feedback & Latency distribution measured over at least 200 sampled follow-ups \\ \hline
NFR-02 & Free-text assessment within 8 s (95th percentile) & Performance & Usability session: explorers find the wait for feedback acceptable & Latency distribution measured over at least 200 sampled assessments \\ \hline
NFR-03 & Progress survives interruption & Reliability & Acceptance test: an explorer closes the browser mid-scenario, returns, and continues & Forced-interruption tests resume at the recorded position \\ \hline
NFR-04 & No progress lost on provider failure; user notified within 10 s & Reliability & Explorers receive a clear message and find their progress intact & Fault injection against the provider interface; zero progress-record loss across injected failures \\ \hline
NFR-05 & Minimal, encrypted, deletable personal data; only de-identified answers sent to AI providers & Security and privacy & Data-protection review against applicable regulation; walkthrough of a deletion request & Data inventory review; encryption check; deletion request removes records; AI requests contain no name, e-mail, or account ID; providers are configured with no training on submitted data and no retention \\ \hline
NFR-06 & Suppress aggregates below a minimum population & Security and privacy & Privacy review confirms small-population figures are withheld and the reason is stated & Aggregates below threshold are withheld \\ \hline
NFR-07 & Disclose AI generation and AI assessment & Transparency & Pilot users can state that scenarios and assessments are machine-produced & Disclosure present at the point of use in every relevant screen \\ \hline
NFR-08 & Vietnamese and English support & Usability & Native speakers review both locales for completeness and accuracy & Both locales render complete content with no missing text \\ \hline
NFR-09 & At least 80\% of readers state the task objective correctly & Usability & Comprehension check with a minimum of 20 representative readers & Validator confirms every scenario states an explicit objective in its context \\ \hline
NFR-10 & New occupations need no code change & Maintainability & A content administrator publishes a new occupation through the authoring interface alone & No code change or redeployment is recorded for the new occupation \\ \hline
NFR-11 & Session cost stays within the configured ceiling & Cost & System administrator reviews cost against the agreed budget & Cost telemetry aggregated per session and compared against the configured ceiling \\ \hline
NFR-12 & Model provider replaceable without content changes & Portability & Domain experts confirm scenario quality is unchanged after a provider switch & Provider substitution test with unchanged content; regression suite passes \\ \hline
NFR-13 & Works in a mobile browser on a slow connection & Compatibility & Explorers complete a scenario on their own phone & Test on reference devices and a throttled network profile agreed for testing \\ \hline
NFR-14 & AI grading validated against expert judgment & Fairness & Domain experts accept the agreement report for their occupation & Cohen's kappa per difficulty stratum recomputed whenever the assessment model changes \\ \hline
NFR-15 & Published scenarios match approved content; AI content during play stays within expert-set limits & Accuracy & Reviewers confirm sampled published scenarios and sampled follow-ups and feedback stay within the approved content & Automated comparison between approved and published content shows no difference; every follow-up and feedback shown passed the output check or is an expert default \\ \hline
NFR-16 & Scenario status changes logged; published content immutable & Auditability & A content administrator traces who approved and published a scenario & Lifecycle log complete for every scenario; write attempts on published content are refused \\ \hline
NFR-17 & Passwords hashed, all traffic encrypted, idle sessions expire, AI prompts protected against injection & Security and privacy & Security review of sign-in, session handling, and AI prompts & Stored passwords are salted hashes; plain-HTTP requests are redirected or refused; sessions end after the configured idle time; explorer answers reach the AI as separated data whose instructions are ignored; an injection test set in the regression suite never changes a score \\ \hline
NFR-18 & Explorer service available within the agreed operating target & Availability & A system administrator reviews monthly availability against the agreed target & Uptime monitoring records availability; planned maintenance is announced in advance \\ \hline
NFR-19 & Explorer screens meet WCAG 2.1 level AA & Accessibility & Users with assistive technology complete a scenario & Automated and manual accessibility audit against WCAG 2.1 AA \\ \hline
\end{longtable}
}

Design constraints are recorded separately from business rules, since they arise from implementation choices rather than from the domain. Each derives from at least one non-functional requirement: the requirement states what must be achieved, and the constraint fixes or prohibits a design choice so that it can be.

{\small
\begin{longtable}{|>{\raggedright\arraybackslash}p{\dimexpr 0.103\textwidth-2\tabcolsep-\arrayrulewidth\relax}|>{\raggedright\arraybackslash}p{\dimexpr 0.458\textwidth-2\tabcolsep-\arrayrulewidth\relax}|>{\raggedright\arraybackslash}p{\dimexpr 0.126\textwidth-2\tabcolsep-\arrayrulewidth\relax}|>{\raggedright\arraybackslash}p{\dimexpr 0.309\textwidth-2\tabcolsep-\arrayrulewidth\relax}|}
\caption{Design constraints.} \\
\hline
\textbf{ID} & \textbf{Design constraint} & \textbf{Derived from} & \textbf{Rationale} \\ \hline
\endfirsthead
\hline
\textbf{ID} & \textbf{Design constraint} & \textbf{Derived from} & \textbf{Rationale} \\ \hline
\endhead
DC-01 & The assessor and generator shall use different model families. & NFR-14 & Reduces correlated failure modes and self-preference bias, thereby improving the independence of automated evaluation (Section 2.5) \\ \hline
DC-02 & The system shall support replacing the underlying model provider without requiring changes to the core application logic. & NFR-12, NFR-04 & Dependence on third-party model providers; portability as funding conditions change \\ \hline
DC-03 & Personal data is processed and stored in compliance with Vietnam's Law on Personal Data Protection and its implementing decree, including the provisions for minors and for cross-border transfer. & NFR-05, NFR-06 & Legal obligation (Sections 2.7, 4.1.4) \\ \hline
DC-05 & The platform is delivered as a web application usable through standard browsers without installation. & NFR-13 & Universal access (Section 4.1.5) \\ \hline
\end{longtable}
}

\subsubsection{Data requirements}

{\small
\begin{longtable}{|>{\raggedright\arraybackslash}p{\dimexpr 0.103\textwidth-2\tabcolsep-\arrayrulewidth\relax}|>{\raggedright\arraybackslash}p{\dimexpr 0.199\textwidth-2\tabcolsep-\arrayrulewidth\relax}|>{\raggedright\arraybackslash}p{\dimexpr 0.204\textwidth-2\tabcolsep-\arrayrulewidth\relax}|>{\raggedright\arraybackslash}p{\dimexpr 0.230\textwidth-2\tabcolsep-\arrayrulewidth\relax}|>{\raggedright\arraybackslash}p{\dimexpr 0.259\textwidth-2\tabcolsep-\arrayrulewidth\relax}|}
\caption{Data requirements with source, business verification, and technical verification.} \\
\hline
\textbf{ID} & \textbf{Data requirement} & \textbf{Source} & \textbf{Business verification} & \textbf{Technical verification} \\ \hline
\endfirsthead
\hline
\textbf{ID} & \textbf{Data requirement} & \textbf{Source} & \textbf{Business verification} & \textbf{Technical verification} \\ \hline
\endhead
DR-01 & Occupational reference data (occupations, tasks, skills) & O*NET and ESCO, imported by the content administrator and adapted for Vietnamese practice by domain experts & Domain experts review the adaptation for Vietnamese practice & Imported records resolve against source identifiers \\ \hline
DR-02 & Competency taxonomy with observable indicators & Derived from DR-01 under expert adaptation & Domain experts review definitions and indicators & Every rubric criterion resolves to a taxonomy entry \\ \hline
DR-03 & Occupational adjacency with overlap measures & Computed from DR-01/DR-02 & Domain experts review computed adjacency for plausibility & Overlap values are reproducible from DR-01/DR-02 \\ \hline
DR-04 & Scenario content: context, activities, option scores, rubric, hints, random events, side quests & Author domain expert with AI assistance & Reviewer sign-off recorded before presentation & Context schema validation; automatic checks pass before submission \\ \hline
DR-05 & Generated scenarios with provenance & System & A content administrator traces any served scenario back to its context and model & Every served scenario carries resolvable provenance \\ \hline
DR-06 & User progress: enrolled levels, competence credit capped per scenario, side-quest points, play history, awards & System & Explorers confirm their progress and awards are shown correctly & Progress reconstructs from stored records; replaying a scenario never adds credit beyond the cap \\ \hline
DR-07 & Assessment records, including hint use and explorer flags & System, with flags from explorers & Domain experts audit a sample of assessment records and review flagged content & Records complete for every assessed response; every flag reaches a domain expert of that occupation \\ \hline
DR-08 & Five-attribute occupation ratings (project-defined) & Explorers & Acceptance test of the rating flow after full completion & Ratings exist only for fully completed occupations \\ \hline
DR-09 & Expert reference set stratified by difficulty & Domain experts & Inter-expert agreement recorded for the reference set & Set covers the defined difficulty strata; every item passed validation \\ \hline
DR-10 & Operational telemetry & System & A system administrator answers cost and error questions from the telemetry & Telemetry present and queryable \\ \hline
DR-11 & Play-test session flag & System & Staff confirm dashboards are unchanged after play-testing & Play-test sessions are absent from explorer-facing aggregates \\ \hline
DR-12 & Quiz answers, RIASEC profile with social index, and AI portrait & Explorers and AI model & Pilot users judge the description relevant and not misleading & Profile recalculates from stored answers; retaking the quiz replaces the profile and portrait \\ \hline
DR-13 & Scenario lifecycle record, including review decisions, comments, and revalidation results & System, from staff actions & A content administrator reconstructs a scenario's history & Every transition in SM-01 is logged; every review and revalidation result is stored with actor and time; at most one active scenario per task \\ \hline
DR-14 & Presentation data and asset library & Content administrator with AI & Reviewers check sampled assets for stereotyping & Every presented asset resolves to a library entry with a bias-check result \\ \hline
DR-15 & Occupation catalogue, including level descriptions and future roles & Content administrator and domain experts & A content administrator reviews the catalogue for completeness & Every published occupation has all required fields and a description for every level \\ \hline
DR-16 & User accounts, roles, assignments, and activity log & System administrator; activity log from the system & A system administrator confirms each occupation's staff and reviews an account's activity & Exactly one content administrator and at least one domain expert per occupation; every account action is logged \\ \hline
DR-17 & Platform settings & System administrator & A system administrator reviews settings against policy & Every change is versioned with actor and time \\ \hline
DR-18 & App feedback records & Explorers & A system administrator reviews recent feedback & Each record holds text, time, app version, and handling status \\ \hline
DR-19 & Scenario review criteria & Content administrator & Domain experts confirm the criteria match how they review & Criteria are versioned; each review links to the version it used \\ \hline
\end{longtable}
}

Personal data within DR-06, DR-07, DR-08, DR-11, DR-12, DR-16, and DR-18 is subject to the policies in Section 4.1.4. DR-09 contains no personal data and is retained as a research asset.

\subsection{Summary}

This chapter specified the proposed system against the domain established in Chapters 1 to 3. The business context identified three operating processes --- exploration, content lifecycle, and evaluation --- and twelve business rules, several of which derive from the evidence reviewed in Chapter 2 rather than from product preference: accessible entry levels, competence-gated progression, taxonomy-derived adjacency, mandatory expert validation, and formative-only assessment. It set out a public-benefit business model in which all content is free to users, operation is funded through grants, acknowledged sponsorship, donations, and an intended partnership with the Ministry of Education and Training, and sponsors are excluded from content so that the candour of occupational representation is protected, together with a five-stage roadmap in which each stage is gated by evidence, and a staged organisation of five operating units under a steering committee, with content governed by occupational domain. The system description identified four user groups --- explorers, domain experts, content administrators, and system administrators --- two research problems in AI applied to education --- grounded generation and the validity of stealth assessment, each stated with the comparison that makes its answer knowledge --- and the principal difficulties, the most consequential being generation fidelity, assessment validity, expert availability, cost sustainability without revenue, and the uneven alignment between screen-based tasks and live work. The requirements analysis produced forty-three prioritised user requirements, from which forty-four functional, nineteen non-functional, and nineteen data requirements were derived, each paired with a business and a technical verification method, together with four design constraints, each derived from a non-functional requirement and recorded separately from the business rules. Chapter 5 models the processes and interactions these requirements imply.
